% Generated draft: requires mathematical and source review.
\documentclass[11pt]{article}
\usepackage[margin=25mm]{geometry}
\usepackage{fontspec}
\setmainfont{Latin Modern Roman}
\usepackage{amsmath,amssymb,mathtools,unicode-math}
\setmathfont{Latin Modern Math}
\usepackage{enumitem,needspace,xurl,hyperref}
\setlength{\parindent}{0pt}
\setlength{\parskip}{5pt}
\setlength{\emergencystretch}{3em}
\newcommand{\N}{\mathbb{N}}
\newcommand{\Z}{\mathbb{Z}}
\newcommand{\Q}{\mathbb{Q}}
\newcommand{\R}{\mathbb{R}}
\newcommand{\C}{\mathbb{C}}
\newcommand{\F}{\mathbb{F}}
\newcommand{\A}{\mathbb{A}}
\newcommand{\PP}{\mathbb{P}}
\newcommand{\E}{\mathbb{E}}
\newcommand{\eps}{\varepsilon}
\newcommand{\dd}{\,\mathrm{d}}
\begin{document}
\section*{UnsolvedMath ID: 700203}
\textbf{Source problem code:} catalogue entry
\subsection{Definitions and mathematical statement}
For a group \(H\), define its derived series by \(H^{(0)}=H\) and \(H^{(n+1)}=[H^{(n)},H^{(n)}]\) for \(n\geq 0\). The group \(H\) is soluble if \(H^{(d)}=\{1\}\) for some integer \(d\geq 0\). A group is locally soluble if each of its finitely generated subgroups is soluble.

The source does not define ``binary soluble'' or ``finite rank.'' For this draft, interpret binary solubility as the requirement that \(\langle x,y\rangle\) is soluble for every \(x,y\in G\), and interpret finite rank as finite Pr\"ufer rank: an abelian group \(A\) has finite rank if there exists an integer \(r_A\geq 0\) such that every finitely generated subgroup of \(A\) can be generated by at most \(r_A\) elements. These interpretations require confirmation against the source; in particular, finite Pr\"ufer rank is distinct from finite torsion-free rank.

Under these interpretations, the question is whether every group \(G\) satisfying
\[
\begin{gathered}
\forall x,y\in G\;\exists d_{x,y}\geq 0:
\quad \langle x,y\rangle^{(d_{x,y})}=\{1\},\\
\forall A\leq G\text{ abelian}\;\exists r_A\geq 0\;
\forall B\leq A\text{ finitely generated}:\quad d(B)\leq r_A
\end{gathered}
\]
is locally soluble. Here all displayed bounds are integers, and \(d(B)\) denotes the least number of generators of \(B\). Equivalently, must
\[
\forall F\subseteq G\text{ finite}\;\exists d_F\geq 0:
\quad \langle F\rangle^{(d_F)}=\{1\}
\]
hold? No uniform bound on the derived lengths or on the ranks of different abelian subgroups is assumed, and no finiteness or finite-generation hypothesis on \(G\) is supplied.
\subsection{Short English statement}
Must a group in which every two elements generate a soluble subgroup, and every abelian subgroup has finite Pr\"ufer rank, be locally soluble? This formulation uses explicit interpretations of two terms left undefined in the supplied source and requires human mathematical review.
\subsection{Sources}
\begin{itemize}
\item Supplied local UnsolvedMath registry record, permanent ID \(700203\), exact source code \texttt{catalogue entry}, titled ``Kourovka Notebook Problem 3.55.'' The record labels the problem open; this status is supplied metadata and has not been independently verified.
\item The supplied attribution is to S. P. Strunkov, with proposed year 1969. The background identifies \emph{The Kourovka Notebook (Unsolved Problems in Group Theory)}, 3rd issue (1969), original problem number 3.55, as reproduced in the 21st issue edition (2026), \texttt{arXiv:1401.0300v48}. This bibliographic attribution is supplied and has not been independently checked.
\item The supplied collection description identifies E. I. Khukhro and V. D. Mazurov as editors of the Kourovka Notebook collection covering Issues 1--21 (1965--2026).
\end{itemize}
\end{document}
